\section{Putting Everything Together}
\label{sec:together}

We compose the building blocks of Section~\ref{sec:blocks}
with the commitment interface of Section~\ref{sec:pcs} to prove
the PLONK relation in Definition~\ref{def:plonk}.
For a single instance, we pack consecutive blocks of each wire
vector and encode copy constraints over
\(\Field\subset\mathbb K\subset\ExtensionField\).
For multiple instances of one circuit, we pack corresponding
wires across instances. The shared selectors and wiring then act
directly on the packed vectors over \(\ExtensionField\).
Sections~\ref{sec:plonk-packing} and~\ref{sec:multi-instance-packing}
take \(\Field=\mathbb F_p\) for an odd prime \(p\).
Section~\ref{sec:binary-packing} gives the binary-field instantiation.

\subsection{Packing for a Single Instance}
\label{sec:plonk-packing}

Algorithm~\ref{alg:single-instance-packing} in
Appendix~\ref{appx:pseudocode} summarizes the protocol.

Let \(N=m\kappa\) be the number of gates, where \(m,\kappa\)
are powers of two and \(p=|\Field|>\kappa\).
The circuit fixes the selectors and augmented wiring permutation
of Definition~\ref{def:plonk}; the public statement supplies
\(\bm{\mathsf{io}}\).
The prover arranges each of
\(\bm w_L,\bm w_R,\bm w_O\) as an \(m\times\kappa\) matrix
and commits using Section~\ref{sec:pcs}.
These three commitments are fixed before any verifier challenge.
The public I/O and the product \(\bm w_M\coloneqq\bm w_L\circ\bm w_R\)
require no additional witness commitments.

\paragraph{Gate constraints.}
The verifier samples
\(\ext{\bm x}\sample\ExtensionField^{\log_2N}\) and sets
\(\ext{\bm\xi}=\tensorspan{\ext{\bm x}}\).
It checks one random projection of the gate residual:
\begin{equation}
    \left\langle \ext{\bm\xi},
    \bm q_L\circ\bm w_L+\bm q_R\circ\bm w_R
    +\bm q_M\circ(\bm w_L\circ\bm w_R)
    +\bm q_O\circ\bm w_O+\bm q_C\right\rangle=0.
\label{eq:plonkWeight}
\end{equation}
The prover sends
\(\ext c_L,\ext c_R,\ext c_M,\ext c_O,\ext c_C\);
the verifier checks their sum and the following claims:
\begin{align}
    \ext c_L+\ext c_R+\ext c_M+\ext c_O+\ext c_C&=0,
    \nonumber\\
    \ext c_L&=\langle\ext{\bm\xi},\bm q_L\circ\bm w_L\rangle,
    \quad
    \ext c_R=\langle\ext{\bm\xi},\bm q_R\circ\bm w_R\rangle,
    \quad
    \ext c_O=\langle\ext{\bm\xi},\bm q_O\circ\bm w_O\rangle,
    \label{eq:plonkWeightedHadamard1}\\
    \ext c_M&=\langle\ext{\bm\xi},\bm q_M\circ\bm w_M\rangle,
    \quad
    \ext c_C=\langle\ext{\bm\xi},\bm q_C\rangle.
    \label{eq:plonkWeightedHadamard2}
\end{align}
The three product claims in Eq.~\eqref{eq:plonkWeightedHadamard1}
are instances of Eq.~\eqref{eq:hadamard-claim}.
Algorithm~\ref{alg:hadamard} reduces each to two packed evaluation
claims, one for a selector and one for a committed wire vector.

Apply the same reduction to the claim for \(\ext c_M\).
Its terminal claim on \(\bm q_M\) is opened directly.
Its terminal claim on the virtual vector \(\bm w_M\) is another
weighted Hadamard-product claim on \(\bm w_L,\bm w_R\), so a
second invocation of Algorithm~\ref{alg:hadamard} authenticates it.
The induced column weights may be arbitrary, while the row weights
remain a tensor product; these are precisely the queries supported
by Eq.~\eqref{eq:hadamard-claim} and Section~\ref{sec:pcs}.
Thus the gate argument uses five packed Hadamard proofs and
ten packed openings: five selector openings and five wire openings,
including the direct opening for \(\ext c_C\) among the selectors.

\paragraph{Joint encoding of positions and permutation images.}
Pad the augmented assignment to a power-of-two length
\(L\geq3N+n_{\mathsf{io}}\), and extend \(\sigma\) by fixed points
on the zero padding.
Positions follow their order in
\((\bm w_L\mathbin{\|}\bm w_R\mathbin{\|}\bm w_O
\mathbin{\|}\bm{\mathsf{io}}\mathbin{\|}\bm0)\).
Let \(\underline p=2^s\) be the largest power of two below \(p\),
set \(\gamma=\lceil\log_{\underline p}L\rceil\), and define
\(\delta_t=2^{s-t}\in\Field\) for \(t\in[s]\).
Since \(\underline p<p\), the bit-to-field map
\((b_1,\ldots,b_s)\mapsto\sum_{t=1}^s\delta_tb_t\)
is injective on \(\{0,1\}^s\).
Choose compatible fields and bases:
\begin{equation}
    \Field\subset\mathbb K\subset\ExtensionField,
    \qquad [\mathbb K:\Field]=d\geq\gamma+1,
    \qquad [\ExtensionField:\mathbb K]
       =\kappa_{\mathbb K}=\kappa/d\geq2,
\label{eq:together-field-tower}
\end{equation}
where \(d\mid\kappa\) and \(\kappa_{\mathbb K}\) is a power of two.
Fix an \(\Field\)-basis
\((\ext\theta_1,\ldots,\ext\theta_d)\) of \(\mathbb K\).
Write \(i-1\) as a most-significant-bit-first string of
\(\gamma s\) bits, with leading zeros as necessary.
Split it into \(\gamma\) consecutive blocks of \(s\) bits,
and let \(e_{i,j}\in\Field\) be the image of block \(j\)
under the bit-to-field map. These radix-\(\underline p\) digits
encode the lexicographically numbered positions injectively. Define:
\begin{equation}
    \ext\tau_i=\sum_{j=1}^{\gamma}e_{i,j}\ext\theta_j,
    \qquad
    \ext\tau_{\sigma,i}
       =\sum_{j=1}^{\gamma}e_{\sigma(i),j}\ext\theta_j.
\label{eq:together-wiring-encoding}
\end{equation}
The setup commits to the \(\gamma\) digit vectors
\(\bm\tau_{\sigma,j}=(e_{\sigma(i),j})_{i=1}^L\in\Field^L\)
using the \(\Field\)-based packed PCS.
The identity labels \(\ext{\bm\tau}\) need no commitment.

\begin{proposition}[Efficient MLE evaluation of position encodings]
\label{proposition:position-encoding-evaluation}
Let \(L=2^\mu\), \(\underline p=2^s\), and
\(\gamma=\lceil\mu/s\rceil\), with digit coefficients
\(\delta_1,\ldots,\delta_s\) as above.
For \(\ext{\bm x}\in\ExtensionField^\mu\), prepend zeros and
split the resulting vector into \(\gamma\) blocks of length \(s\):
\[
    \widehat{\bm x}
       =\bm0_{\gamma s-\mu}\mathbin{\|}\ext{\bm x}
       =\widehat{\bm x}^{(1)}\mathbin{\|}\cdots
          \mathbin{\|}\widehat{\bm x}^{(\gamma)}.
\]
Then the position-label MLE is:
\begin{equation}
    \left\langle\ext{\bm\tau},
       \tensorspan{\ext{\bm x}}\right\rangle
       =\sum_{j=1}^{\gamma}\ext\theta_j
          \sum_{t=1}^{s}\delta_t\widehat{x}^{(j)}_t.
\label{eq:position-encoding-evaluation}
\end{equation}
It can be evaluated in \(O(\log L)\) field operations.
The same bound holds for a query given by arbitrary two-entry
tensor factors.
\end{proposition}

The proof is deferred to
Appendix~\ref{appx:proof-position-encoding}.

\paragraph{Copy constraints.}
Reserve the last basis coordinate for the wire value:
\begin{equation}
    \ext{\bm a}=\ext\theta_d\bm w+\ext{\bm\tau},
    \qquad
    \ext{\bm b}=\ext\theta_d\bm w+\ext{\bm\tau}_\sigma
       \in\mathbb K^L.
\label{eq:together-copy-operands}
\end{equation}
The labels are distinct, and their span does not contain
\(\ext\theta_d\). Hence joint encoding is injective and gives:
\begin{equation}
    \sigma(\bm w)=\bm w
    \quad\Longleftrightarrow\quad
    \ext{\bm a}\stackrel{\mathsf{perm}}{=}\ext{\bm b}.
\label{eq:together-copy-equivalence}
\end{equation}
Treating \(\mathbb K\) as the base field, apply
Algorithm~\ref{alg:permutation} with length \(L\), packing width
\(\kappa_{\mathbb K}\), and extension field \(\ExtensionField\).
The two operands are determined by the existing commitments,
so they need not be committed separately.

Each terminal copy query splits into evaluations of
\(\bm w_L,\bm w_R,\bm w_O\), the public I/O, and the labels.
The three wire segments are aligned power-of-two blocks, so their
restricted queries retain the packed tensor form.
The verifier evaluates the public I/O and identity labels directly.
Only the query to \(\ext{\bm b}\) requires the \(\gamma\)
permuted-digit openings.
Together with the gate argument, the setup tables contain
\(5N+\gamma L\) base-field entries, and the online protocol uses
\(16+\gamma\) packed openings. No \(\mathbb K\)-based PCS is needed.

\paragraph{Efficiency and soundness.}
Let \(\epsilon_{\mathsf{HP}}\) be the bound in
Proposition~\ref{proposition:hadamard-complexity} for length \(N\),
and let \(\epsilon_{\mathsf{GP}}^{\mathbb K}\) be the one-grand-product
bound in Proposition~\ref{proposition:permutation-complexity}
with parameters
\((\mathbb K,\ExtensionField,L,L/\kappa_{\mathbb K},
\kappa_{\mathbb K})\).
These bounds exclude terminal openings.
Let \(\epsilon_{\mathsf{PCS}}\) bound each underlying adaptive
evaluation error, and let \(\epsilon_{\mathsf{bind}}\) bound the
joint binding failure of the three witness commitments.
Write the packed opening error for a length-\(n\) base-field
vector as:
\begin{equation}
    \epsilon_{\mathsf{open}}(n)
       =\epsilon_{\mathsf{PCS}}
        +\frac{1+2\log_2(n/\kappa)}{|\ExtensionField|}
        +\left(\frac{2\log_2(n/\kappa)+2}{p}\right)^k.
\label{eq:together-opening-error}
\end{equation}
Use a common repetition count \(k\) from
Eq.~\eqref{eq:pcs-repetition-number} for the largest row count
\(L/\kappa\), assuming \(p>2\log_2(L/\kappa)+2\).
For efficiency, use schoolbook arithmetic in compatible polynomial
bases for the tower and the corresponding product basis for
the \(\Field\)-based PCS.

\begin{theorem}[Single-instance packing]
\label{theorem:single-instance-packing}
Under the parameters above, assume \(m\geq\kappa\) and the PCS
assumptions of Proposition~\ref{proposition:packed-pcs}.
The single-instance protocol is complete and has soundness error
at most:
\begin{equation}
\begin{aligned}
    &\epsilon_{\mathsf{bind}}
       +\frac{\log_2N}{|\ExtensionField|}
       +5\epsilon_{\mathsf{HP}}
       +\min\left\{1,\frac{L-1}{|\ExtensionField|}\right\}\\
    &\qquad{}
       +2\epsilon_{\mathsf{GP}}^{\mathbb K}
       +16\epsilon_{\mathsf{open}}(N)
       +\gamma\epsilon_{\mathsf{open}}(L).
\end{aligned}
\label{eq:together-soundness}
\end{equation}
With preprocessed circuit tables and packing polynomials, the
online costs, including the packed opening reductions, are:
\begin{equation}
\begin{aligned}
    \text{Prover time}\quad
       &O\left((\kappa+k)N+
          [\kappa d+\gamma(\kappa+k)]L
          +\kappa^2(\log_2L)^2\right)\bb,\\
    \text{Verifier time}\quad
       &O\left(\kappa^2\left[
          (\gamma+1)\kappa+(\log_2L)^2
          +(\gamma+1)k\log_2L+n_{\mathsf{io}}
          \right]\right)\bb,\\
    \text{Communication}\quad
       &O\left((\gamma+1)\kappa^2+\kappa(\log_2L)^2
          +(\gamma+1)k\log_2L\right),\\
    \text{Prover space}\quad
       &O(N+Ld).
\end{aligned}
\label{eq:together-single-efficiency}
\end{equation}
Communication and space are measured in \(\Field\)-coordinates.
The bounds exclude one-time circuit preprocessing and the
underlying PCS costs. Total online costs additionally include
three underlying witness commitments and \(16+\gamma\)
underlying openings: sixteen at vector length \(m\) and
\(\gamma\) at vector length \(L/\kappa\) over \(\ExtensionField\).
The underlying PCS workspace is added to the space bound.
\end{theorem}

The proof is deferred to Appendix~\ref{appx:proof-single-instance}.

\subsection{Packing Multiple Instances of the Same Circuit}
\label{sec:multi-instance-packing}

Algorithm~\ref{alg:multiple-instance-packing} in
Appendix~\ref{appx:pseudocode} summarizes the protocol.

Consider \(\kappa\) instances of an \(m\)-gate circuit, with
\(N=m\kappa\), and arrange their wire assignments columnwise:
\[
    \bm W_\ell
       =\concat{\bm w_\ell^{(1)},\ldots,\bm w_\ell^{(\kappa)}}
       \in\Field^{m\times\kappa},
       \qquad \ell\in\{L,R,O\}.
\]
Fix the basis \(\ext{\bm\beta}\) from Section~\ref{sec:pcs}.
The prover commits to the three packed vectors; the product
vector remains virtual:
\begin{equation}
    \ext{\bm w}_\ell=\bm W_\ell\ext{\bm\beta}
       \quad(\ell\in\{L,R,O\}),
    \qquad
    \ext{\bm w}_M=(\bm W_L\circ\bm W_R)\ext{\bm\beta}.
\label{eq:together-multiple-wires}
\end{equation}
Each \(\ExtensionField\)-vector uniquely determines its
base-field matrix through the fixed basis.
All commitments precede the verifier's challenges.

\paragraph{Gate constraints.}
The selectors now have length \(m\) and are shared by all
instances. Multiplying the matrix of gate residuals by
\(\ext{\bm\beta}\) yields:
\begin{equation}
    \bm q_L\circ\ext{\bm w}_L
    +\bm q_R\circ\ext{\bm w}_R
    +\bm q_M\circ\ext{\bm w}_M
    +\bm q_O\circ\ext{\bm w}_O
    +\left(\sum_{j=1}^{\kappa}\ext\beta_j\right)\bm q_C
       =\bm0_m.
\label{eq:together-multiple-gates}
\end{equation}
This identity is equivalent to the gate constraints of all
\(\kappa\) instances, by linear independence of the basis.
In particular, the constant-selector term is packed along
with the witness columns.

The verifier samples a random tensor projection of the length-\(m\)
residual. A degree-three sum-check over \(\ExtensionField\)
reduces it to five selector evaluations and four packed wire
evaluations at one common point.
The five selectors and
\(\ext{\bm w}_L,\ext{\bm w}_R,\ext{\bm w}_O\)
are opened directly using the underlying multilinear PCS.
For the remaining evaluation, with tensor weight
\(\ext{\bm\chi}\), the parties prove:
\begin{equation}
    \left\langle\ext{\bm w}_M,\ext{\bm\chi}\right\rangle
       =\ext{\bm\chi}^{\top}
          (\bm W_L\circ\bm W_R)\ext{\bm\beta}.
\label{eq:together-multiple-product-opening}
\end{equation}
Algorithm~\ref{alg:hadamard} applies with column weights
\(\ext{\bm\beta}\) and row weights \(\ext{\bm\chi}\).
Its two terminal queries are authenticated against the existing
wire commitments by the packed opening protocol.
Thus only one packed Hadamard proof is needed for the entire batch.

\paragraph{Copy constraints.}
Here \(L\geq3m+n_{\mathsf{io}}\) is a power of two bounding
the augmented assignment length of \emph{one} instance.
Pack the public-I/O columns in the same basis and concatenate:
\[
    \ext{\bm w}
       =\ext{\bm w}_L\mathbin{\|}\ext{\bm w}_R
          \mathbin{\|}\ext{\bm w}_O
          \mathbin{\|}\ext{\bm{\mathsf{io}}}\mathbin{\|}\bm0
       \in\ExtensionField^L.
\]
The shared permutation acts on rows, so
\(\sigma(\ext{\bm w})=\ext{\bm w}\) holds if and only if it
holds in every instance.
Use the same radix-\(\underline p\) position digits, now with
\(\gamma=\lceil\log_{\underline p}L\rceil\leq\kappa\), and embed
the labels directly in \(\ExtensionField\):
\begin{equation}
    \ext\tau_i=\sum_{j=1}^{\gamma}e_{i,j}\ext\beta_j,
    \qquad
    \ext\tau_{\sigma,i}
       =\sum_{j=1}^{\gamma}e_{\sigma(i),j}\ext\beta_j.
\label{eq:together-multiple-labels}
\end{equation}
The setup commits to \(\ext{\bm\tau}_\sigma\in\ExtensionField^L\)
using the underlying PCS. Proposition~\ref{proposition:position-encoding-evaluation}
evaluates the identity labels with \(\ext\theta_j=\ext\beta_j\).

After the witness commitments are fixed, sample independent
\(\ext y,\ext z\sample\ExtensionField\).
The parties prove equality of the following grand products:
\begin{equation}
    \prod_{i=1}^{L}
       (\ext z+\ext w_i+\ext y\ext\tau_i)
       =
    \prod_{i=1}^{L}
       (\ext z+\ext w_i+\ext y\ext\tau_{\sigma,i}).
\label{eq:together-multiple-copy}
\end{equation}
Unlike the single-instance encoding, the values occupy all
\(\kappa\) basis coordinates. The random coefficient \(\ext y\)
separates them from the position labels without an intermediate
field \(\mathbb K\).

Each grand product is proved by a binary product tree directly
over \(\ExtensionField\). At each level, a degree-three sum-check
reduces the current weighted product claim to two child
evaluations. A fresh random linear combination merges them into
one evaluation of the preceding layer.
The two trees terminate in one evaluation of each virtual operand.
These require six direct wire openings and one direct opening of
\(\ext{\bm\tau}_\sigma\); the identity labels and public I/O
are evaluated by the verifier.
The circuit preprocessing uses five length-\(m\) selector
commitments and one length-\(L\) label commitment over
\(\ExtensionField\), shared across the batch.

\paragraph{Efficiency and soundness.}
Use \(\epsilon_{\mathsf{HP}}\), \(\epsilon_{\mathsf{bind}}\),
and \(\epsilon_{\mathsf{PCS}}\) as above.
In this subsection, choose \(k\) from
Eq.~\eqref{eq:pcs-repetition-number} with \(m\) rows,
assuming \(p>2\log_2m+2\), and use
\(\epsilon_{\mathsf{open}}(N)\) from
Eq.~\eqref{eq:together-opening-error}.
Set \(\mu=\log_2L\).
The two random coefficients in Eq.~\eqref{eq:together-multiple-copy}
are analyzed jointly, so no separate position-collision term is needed.

\begin{theorem}[Multiple-instance packing]
\label{theorem:multiple-instance-packing}
Let \(m,\kappa\) be powers of two with \(m\geq\kappa\geq2\),
\(N=m\kappa\), and \(p>\kappa\).
Under the parameters and PCS assumptions above, the protocol
for \(\kappa\) instances of the same circuit is complete and has
soundness error at most:
\begin{equation}
\begin{aligned}
    &\epsilon_{\mathsf{bind}}
       +\frac{4\log_2m}{|\ExtensionField|}
       +\epsilon_{\mathsf{HP}}
       +\min\left\{1,\frac{L-1}{|\ExtensionField|}\right\}\\
    &\qquad{}
       +\frac{3\mu(\mu-1)+2\mu}{|\ExtensionField|}
       +15\epsilon_{\mathsf{PCS}}
       +2\epsilon_{\mathsf{open}}(N).
\end{aligned}
\label{eq:together-multiple-soundness}
\end{equation}
With preprocessed circuit tables and packing polynomials, the
online costs under schoolbook arithmetic are:
\begin{equation}
\begin{aligned}
    \text{Prover time}\quad
       &O\left((\kappa+k)N+
          \kappa^2[L+(\log_2L)^2]\right)\bb,\\
    \text{Verifier time}\quad
       &O\left(\kappa^2\left[
          \kappa+(\log_2L)^2+k\log_2m+n_{\mathsf{io}}
          \right]\right)\bb,\\
    \text{Communication}\quad
       &O\left(\kappa^2+\kappa(\log_2L)^2
          +k\log_2m\right),\\
    \text{Prover space}\quad
       &O(N+\kappa L).
\end{aligned}
\label{eq:together-multiple-efficiency}
\end{equation}
Communication and space are measured in \(\Field\)-coordinates.
The bounds include the two packed opening reductions but exclude
one-time preprocessing and the underlying PCS costs.
There are three online witness commitments and seventeen
underlying openings: sixteen at length \(m\) and one at length
\(L\) over \(\ExtensionField\).
The underlying PCS workspace is added to the space bound.
In particular, if \(n_{\mathsf{io}}=O(m)\) and \(L=O(m)\),
the prover cost excluding the underlying PCS is
\(O((\kappa+k)N+\kappa^2(\log_2m)^2)\bb\) for the whole batch,
or \(O((\kappa+k)m+\kappa(\log_2m)^2)\bb\) per instance.
\end{theorem}

The proof is deferred to
Appendix~\ref{appx:proof-multiple-instance}.

Both compositions establish arguments for the constraint
relation. Zero knowledge additionally requires compatible masking
of the subprotocol transcripts.

\subsection{Binary-Field Instantiation}
\label{sec:binary-packing}

The reductions above also support binary circuits by using a
small binary extension as the PIOP working field. Choose the tower:
\[
    \mathbb F_2\subset\mathbb K=\mathbb F_{2^8}
       \subset\ExtensionField,
    \qquad [\mathbb K:\mathbb F_2]=8,
    \qquad [\ExtensionField:\mathbb K]=\kappa_{\mathbb K}.
\]
Arithmetize the binary circuit over \(\mathbb K\), including
\(w^2-w=0\) for the intended binary wires before defining
\(m,N\). Instantiate the building blocks, packed PCS, and
multiple-instance composition with
\((\Field,p,\kappa)=(\mathbb K,256,\kappa_{\mathbb K})\).
Here \(N=m\kappa_{\mathbb K}\) counts working-field entries.

\paragraph{Packing and openings.}
Choose \(\kappa_{\mathbb K}\) distinct interpolation nodes in
\(\mathbb K\) and represent \(\ExtensionField\) in a
\(\mathbb K\)-basis. All coordinate matrices in
Section~\ref{sec:pcs} use this basis, so the repeated query-side
checks and their terminal coordinate-recovery identities are
\(\mathbb K\)-linear. The gate and copy challenges remain in
\(\ExtensionField\). This changes the working field of the
whole reduction, including its coordinates, rather than only
the challenge distribution.

\paragraph{Position encoding.}
Set \(s=8\), \(\underline p=256\), and
\(\gamma=\lceil\log_{256}L\rceil\leq\kappa_{\mathbb K}\).
Choose an \(\mathbb F_2\)-basis
\((\delta_1,\ldots,\delta_8)\) of \(\mathbb K\) and encode each
eight-bit position block as \(\sum_{t=1}^8\delta_tb_t\).
The resulting digits are injective, and the labels embed directly
in \(\ExtensionField\) by Eq.~\eqref{eq:together-multiple-labels}.
The MLE identity in Eq.~\eqref{eq:position-encoding-evaluation}
holds for these coefficients by the same multilinear interpolation
argument in Appendix~\ref{appx:proof-position-encoding}.
The randomized copy argument therefore needs no further field.
Here \(\mathbb K\) is a working field, not the joint
value-and-position field of Section~\ref{sec:plonk-packing}:
using \(\mathbb F_{2^8}\) for the latter over \(\mathbb F_2\)
would require \(1+\lceil\log_2L\rceil\leq8\), hence \(L\leq2^7\).

\paragraph{Using a binary linear code.}
Section~\ref{sec:pcs-base-code} gives the code-switching
instantiation of the PCS. A binary generator matrix extends
first to \(\mathbb K\) and then to \(\ExtensionField\) without
changing its entries. The source encoding can therefore be
computed on the eight bit-coordinate matrices over \(\mathbb F_2\),
while the packed opening reduction uses the
\(\ExtensionField/\mathbb K\) basis.
Code switching supplies the terminal extension-field opening
through a characteristic-two-compatible BaseFold
instantiation~\cite{eurocrypt/DiamondP26}.

\paragraph{Efficiency and soundness.}
When using the packed PCS of Section~\ref{sec:pcs}, the proof of
Theorem~\ref{theorem:multiple-instance-packing} is unchanged:
it uses field identities, injective encodings, and polynomial
degree bounds, all valid in characteristic two.
Its efficiency and soundness bounds thus hold with
\(p=256\) and \(\kappa=\kappa_{\mathbb K}\), subject to the
stated field-size and packing conditions.
In particular, the repeated-check error term is
\(((2\log_2m+2)/256)^k\).
Arithmetic costs are measured in \(\mathbb K\)-operations,
and communication and space in \(\mathbb K\)-coordinates.
The working-field PIOP steps therefore use byte-field arithmetic;
the packing across instances is still performed through a
\(\mathbb K\)-basis of \(\ExtensionField\).

For the code-based instantiation, use
the underlying PCS error and costs of the chosen Blaze
instantiation, as specified in Section~\ref{sec:pcs-base-code}.
The working-field packed-opening reduction and the gate,
Hadamard, and copy-reduction bounds remain unchanged.
